%============================================================
% CH1 · LESSON 4 — SCANNING
%============================================================
\subsection{Scanning: One Pass, Two Passes, and Extra Arrays}
\label{ssec:ch1-scanning}

A \textbf{scan} visits the elements in order, $a_0, a_1, \dots, a_{n-1}$, and updates a small amount of \emph{state} at each step. It is the most common loop in this book and, thanks to the cache (Lesson~\ref{ssec:ch1-array}), the fastest. The design question is always the same: \emph{what must I remember about the part already read in order to finish the job?} This lesson answers it for the three typical situations --- a constant amount of state, a second pass, and a whole second array --- and ends with the three kinds of ``pieces'' of an array that problems ask about.

%------------------------------------------------------------
\subsubsection*{1. One pass with running state}

\begin{diagrambox}
\begin{tikzpicture}
    \dpindex{oi}{-0.05}{$i$}{0}{6}
    \dparray{o}{0.6}{\texttt{a}}{3/adone, 7/adone, 2/adone, 7/aread, 5/acell, 7/acell, 1/acell}
    \aptrup{o3}{i}{accentorange}
    \arange{oi0}{oi2}{read: $[0, i)$}{tabulationblue}
    \node[rectangle, rounded corners=3pt, draw=successgreen, fill=solutionbg, font=\scriptsize, align=left,
          anchor=west, text width=4.4cm] at (7,0.6)
        {\textbf{state after $[0,3)$}\\ \texttt{sum} $= 12$, \texttt{mx} $= 7$, \texttt{cnt} $= 1$\\
         \textbf{update with} $a_3 = 7$:\\ \texttt{sum} $= 19$, \texttt{mx} $= 7$, \texttt{cnt} $= 2$};
\end{tikzpicture}
\figcap{A scan keeps a summary of $[0, i)$ and folds in $a_i$. Here the summary is the sum, the maximum and how often the maximum occurred.}
\end{diagrambox}

Every one-pass algorithm has the same four parts:
\begin{table}[H]
\centering\small
\renewcommand{\arraystretch}{1.2}
\begin{tabularx}{\linewidth}{@{}>{\bfseries}l X@{}}
\toprule
State & the variables that summarise the prefix $[0, i)$ (a sum, a minimum, the two largest values, the length of the current streak, \ldots) \\
Initial value & the summary of the \emph{empty} prefix: $0$ for a sum, $-\infty$ for a maximum, ``nothing seen'' flags \\
Update & how the summary of $[0, i)$ plus the element $a_i$ gives the summary of $[0, i+1)$ \\
Answer & read off the final state (sometimes updated \emph{during} the scan: ``the best value seen at any time'') \\
\bottomrule
\end{tabularx}
\end{table}

The loop invariant is always ``the state is the summary of $[0, i)$'' (Lesson~\ref{ssec:ch1-indices}), and the whole difficulty is in the \textbf{update}: it must be computable from the old state and $a_i$ alone. A sum, a maximum, a count, the two largest values --- yes. The \emph{median} of the prefix --- no: knowing the old median and $a_i$ is not enough to know the new one. Pattern~1 (Section~\ref{sec:ch1-p1}) is the catalogue of states that work.

%------------------------------------------------------------
\subsubsection*{2. Two passes: when the first pass decides what the second looks for}

``How many times does the maximum occur?'' Two natural solutions:
\begin{itemize}[leftmargin=*]
    \item \textbf{Two passes.} Pass 1 finds the maximum $M$. Pass 2 counts the elements equal to $M$. Each pass is trivial and obviously correct.
    \item \textbf{One pass.} Keep the maximum so far and its count; when a \emph{larger} value appears, the old count becomes worthless --- reset it to 1.
\end{itemize}

\begin{lstlisting}[style=cpp]
#include <bits/stdc++.h>
using namespace std;

int main() {
    vector<int> a = {3, 7, 2, 7, 5, 7, 1};
    int n = a.size();
    // two passes: find the maximum, then count it
    int mx = a[0];
    for (int x : a) mx = max(mx, x);
    int cnt2 = 0;
    for (int x : a) if (x == mx) cnt2++;
    // one pass: reset the count whenever a new maximum appears
    int best = a[0], cnt1 = 0;
    for (int x : a) {
        if (x > best) { best = x; cnt1 = 1; }
        else if (x == best) cnt1++;
    }
    cout << "max " << mx << " appears " << cnt2 << " times";
    cout << " (one pass: " << best << ", " << cnt1 << ")\n";
    // prefix maximum and suffix maximum (second arrays)
    vector<int> pre(n), suf(n);
    for (int i = 0; i < n; i++) pre[i] = (i == 0) ? a[i] : max(pre[i - 1], a[i]);
    for (int i = n - 1; i >= 0; i--) suf[i] = (i == n - 1) ? a[i] : max(suf[i + 1], a[i]);
    for (int i = 0; i < n; i++) cout << pre[i] << (i + 1 < n ? ' ' : '\n');   // 3 7 7 7 7 7 7
    for (int i = 0; i < n; i++) cout << suf[i] << (i + 1 < n ? ' ' : '\n');   // 7 7 7 7 7 7 1
}
\end{lstlisting}

Both print ``max 7 appears 3 times''. The one-pass version is not faster in any way that matters ($2n$ versus $n$ steps, both $\bigO{n}$); it is useful when the data can be read only once (a stream). \textbf{Prefer the version you can prove at a glance}, and use two passes freely: the first pass computes a \emph{global} quantity (maximum, total sum, count of zeros, \ldots) that the second pass needs at every position. Problem~\ref{ssec:ch1-advantage} is exactly this: one pass finds the two largest strengths, the second pass answers every participant.

\begin{diagrambox}
\begin{tikzpicture}
    \dparray{p}{0.75}{pass 1}{3/adone, 7/adone, 2/adone, 7/adone, 5/adone, 7/adone, 1/adone}
    \node[anscell, anchor=west] at (6.7,0.75) {$M = 7$};
    \dparray{q}{-0.35}{pass 2}{3/acell, 7/akeep, 2/acell, 7/akeep, 5/acell, 7/akeep, 1/acell}
    \node[anscell, anchor=west] at (6.7,-0.35) {count $= 3$};
    \draw[dep] (7.3,0.4) -- (7.3,0.0);
\end{tikzpicture}
\figcap{Two passes: the first computes a global value, the second uses it at every position (green: equal to $M$).}
\end{diagrambox}

%------------------------------------------------------------
\subsubsection*{3. Prefix and suffix arrays: a preview}

Sometimes the answer for \emph{each} position $i$ depends on everything to its left, or to its right. Then a single number of state is not enough --- we store the running state for every prefix in an array:
\[
\mathit{pre}[i] = \max(a_0, \dots, a_i), \qquad \mathit{suf}[i] = \max(a_i, \dots, a_{n-1}).
\]
Each is one scan (left to right, right to left), and afterwards ``the largest element left of $i$'' and ``the largest element right of $i$'' are $\bigO{1}$ lookups.

\begin{diagrambox}
\begin{tikzpicture}
    \dpindex{fi}{1.35}{$i$}{0}{6}
    \dparray{f}{0.6}{\texttt{a}}{3/acell, 7/acell, 2/acell, 7/acell, 5/acell, 7/acell, 1/acell}
    \dparray{fp}{-0.25}{\textit{pre}}{3/compcell, 7/compcell, 7/compcell, 7/compcell, 7/compcell, 7/compcell, 7/compcell}
    \dparray{fs}{-1.1}{\textit{suf}}{7/basecell, 7/basecell, 7/basecell, 7/basecell, 7/basecell, 7/basecell, 1/basecell}
    \node[note, anchor=west] at (6.6,-0.25) {left $\to$ right};
    \node[note, anchor=west] at (6.6,-1.1) {right $\to$ left};
\end{tikzpicture}
\figcap{Prefix and suffix maxima of $a = (3,7,2,7,5,7,1)$, the values printed by the program above.}
\end{diagrambox}

The same idea with \emph{sums} instead of maxima --- $P[i] = a_0 + \dots + a_{i-1}$, so that the sum of $[l, r)$ is $P[r] - P[l]$ --- is so important that it has its own chapter: Chapter~\ref{ch:prefix_sums}. Here we only need the habit: \emph{if each position needs a summary of one side, precompute it for all positions in one pass.}

%------------------------------------------------------------
\subsubsection*{4. When is a second array really needed?}

\begin{table}[H]
\centering\small
\renewcommand{\arraystretch}{1.25}
\begin{tabularx}{\linewidth}{@{}>{\raggedright\arraybackslash}p{5.4cm} >{\raggedright\arraybackslash}X@{}}
\toprule
\textbf{situation} & \textbf{memory} \\
\midrule
one number for the whole array (sum, max, count, ``is it sorted?'') & $\bigO{1}$: one pass \\
a global value, then an answer per position (max of the others, distance to the total) & $\bigO{1}$ extra: two passes (Problem~\ref{ssec:ch1-advantage}) \\
each position needs a summary of its \emph{left} side only & $\bigO{1}$: the running state at the moment we reach $i$ \emph{is} that summary \\
each position needs its left side \emph{and} its right side & one suffix array (computed first, right to left), the left side as a running variable \\
the output itself is an array & the output array --- or overwrite the input (Lesson~\ref{ssec:ch1-inplace}) \\
the answer depends on values far apart in both directions in a complicated way & sort, hash, or a later chapter \\
\bottomrule
\end{tabularx}
\end{table}

\begin{insightbox}[Left side for free]
A left-to-right scan already ``knows'' the prefix when it stands at $i$. So for a quantity that combines a left summary with $a_i$ --- ``the largest $a_j - a_i$ with $j > i$'', ``the smallest element before $i$'' --- one running variable replaces the whole \textit{pre} array. Only questions that need the \emph{right} side as well force a stored array (or a scan in the other direction).
\end{insightbox}

%------------------------------------------------------------
\subsubsection*{5. Subarrays, subsequences and subsets}

Problems talk about three kinds of ``parts'' of an array, and confusing them changes the answer completely.

\begin{definitionbox}[Three kinds of parts]
Let $a = (a_0, \dots, a_{n-1})$.
\begin{itemize}[leftmargin=*]
    \item A \textbf{subarray} (also: segment, substring for strings) is a contiguous block $a_l, a_{l+1}, \dots, a_{r-1}$ --- a half-open range $[l, r)$ with $l < r$.
    \item A \textbf{subsequence} keeps some of the elements \emph{in their original order}, not necessarily adjacent: $a_{i_1}, a_{i_2}, \dots, a_{i_k}$ with $i_1 < i_2 < \dots < i_k$.
    \item A \textbf{subset} is a choice of elements where the order does not matter (often the problem allows you to rearrange them anyway).
\end{itemize}
\end{definitionbox}

\begin{diagrambox}
\begin{tikzpicture}
    \begin{scope}
        \dparray{u}{0}{subarray}{3/acell, 1/awin, 4/awin, 1/awin, 5/acell, 9/acell}
        \arange{u1}{u3}{$[1,4)$: $(1,4,1)$}{memoorange}
    \end{scope}
    \begin{scope}[yshift=-1.95cm]
        \dparray{v}{0}{subsequence}{3/akeep, 1/acell, 4/akeep, 1/acell, 5/acell, 9/akeep}
        \draw[dep] (v0.north) to[bend left=30] (v2.north);
        \draw[dep] (v2.north) to[bend left=30] (v5.north);
        \node[note, anchor=west] at (5.6,0) {$(3, 4, 9)$: indices $0 < 2 < 5$};
    \end{scope}
    \begin{scope}[yshift=-3.3cm]
        \dparray{w}{0}{subset}{3/akeep, 1/acell, 4/akeep, 1/acell, 5/acell, 9/akeep}
        \node[note, anchor=west] at (5.6,0) {$\{9, 3, 4\}$ = $\{3, 4, 9\}$: order irrelevant};
    \end{scope}
\end{tikzpicture}
\figcap{The same array, three kinds of parts. Every subarray is a subsequence; every subsequence determines a subset.}
\end{diagrambox}

\textbf{How many are there?} This decides whether ``try all of them'' is an option.
\begin{itemize}[leftmargin=*]
    \item \textbf{Subarrays:} choose $l < r$ from the $n+1$ boundary positions $0, \dots, n$: $\binom{n+1}{2} = \dfrac{n(n+1)}{2}$ non-empty subarrays --- $\bigO{n^2}$. Enumerating them all is fine for $n \le 5000$ if each is handled in $\bigO{1}$ (extend $r$ with running state), and fine for $n \le 500$ even with an $\bigO{n}$ inner loop.
    \item \textbf{Subsequences:} every element is either taken or not: $2^n$ choices, $2^n - 1$ non-empty --- exponential. Trying them all works only for $n \le 20$--$25$; for larger $n$ we need structure (greedy, or DP --- the \textit{Dynamic Programming} book).
    \item \textbf{Subsets} of $n$ elements: also $2^n$ (as many as subsequences, since a subset of positions is a subsequence). If the array has repeated values and only the multiset of values counts, there are fewer distinct ones.
\end{itemize}

\begin{lstlisting}[style=cpp]
#include <bits/stdc++.h>
using namespace std;

int main() {
    vector<char> a = {'a', 'b', 'c', 'd'};
    int n = a.size();
    int subarrays = 0;
    for (int l = 0; l < n; l++)                // every subarray is a half-open [l, r)
        for (int r = l + 1; r <= n; r++) subarrays++;
    int subsequences = 0;
    for (int mask = 1; mask < (1 << n); mask++) { // bit i set = a[i] chosen, order kept
        string s;
        for (int i = 0; i < n; i++) if (mask >> i & 1) s += a[i];
        if (n == 4 && s.size() == 2) cout << s << " ";
        subsequences++;
    }
    cout << "\nn = " << n << ": " << subarrays << " non-empty subarrays (n(n+1)/2 = "
         << n * (n + 1) / 2 << "), " << subsequences << " non-empty subsequences (2^n - 1)\n";
}
\end{lstlisting}

For \texttt{abcd} it prints the six subsequences of length 2 (\texttt{ab ac bc ad bd cd}; only \texttt{ab}, \texttt{bc}, \texttt{cd} are also subarrays) and the counts $10$ and $15$.

\begin{table}[H]
\centering\small
\renewcommand{\arraystretch}{1.2}
\begin{tabular}{@{}lccccl@{}}
\toprule
 & $n = 4$ & $n = 20$ & $n = 10^3$ & $n = 2\cdot10^5$ & enumerate all? \\
\midrule
subarrays $n(n+1)/2$ & 10 & 210 & 500500 & $\approx 2\cdot10^{10}$ & up to $n \approx 5000$ \\
subsequences $2^n - 1$ & 15 & 1048575 & --- & --- & up to $n \approx 20$ \\
\bottomrule
\end{tabular}
\end{table}

\begin{warningbox}[Read the word carefully]
``Subarray'', ``contiguous subsequence'', ``segment'' and ``substring'' mean contiguous. ``Subsequence'' means order kept, gaps allowed. ``Choose some elements'' or ``subset'' usually means order is irrelevant. The maximum sum of a \emph{subsequence} is the sum of all positive elements (take every positive one); the maximum sum of a \emph{subarray} needs Kadane's algorithm (\dpbook{1}{Pattern 5}). Same array, same word ``sum'', different problems.
\end{warningbox}

\begin{checklistbox}[Lesson summary]
\begin{itemize}[leftmargin=*]
    \item[$\square$] Design a scan by its state: initial value (empty prefix), update ($[0,i) + a_i$), answer.
    \item[$\square$] A global value needed at every position $\Rightarrow$ two passes; it is still $\bigO{n}$.
    \item[$\square$] Left side: a running variable. Right side too: a suffix array (or a reverse scan).
    \item[$\square$] Subarrays $n(n+1)/2$ (contiguous), subsequences $2^n - 1$ (order kept), subsets $2^n$ (order ignored).
\end{itemize}
\end{checklistbox}
